\section{Methodik: Neighbor Context Detection im P2P-Netzwerk}
\label{sec:methodology}

\subsection{Inhaltliche Stichpunkte}
\begin{itemize}
  \item \textbf{Lokale Baseline}
  \begin{itemize}
    \item Random-Forest-Klassifikator auf zusammengeführten Trainingsdaten
          der Knoten A bis H.
    \item Eingaben: shunt\_voltage, bus\_voltage\_V, current\_mA,
          power\_mW und State; Zielklassen Backdoor, none und syn-flood.
    \item Lokale Modellwahrscheinlichkeiten für die Angriffsklassen werden
          als Grundlage für Peer-Berichte gespeichert.
  \end{itemize}
  \item \textbf{Kontextaggregation}
  \begin{itemize}
    \item Pro Zeitschritt aggregiert jeder Zielknoten Berichte der sieben
          übrigen Knoten; der Zielknoten selbst wird nicht als Peer gezählt.
    \item Die Peer-Berichte werden anhand einer Sigmoid-Funktion gewichtet:
          Basisgewicht 1, Maximalgewicht 3, Steilheit 10 und Wendepunkt bei
          einer Angriffswahrscheinlichkeit von 0{,}5.
    \item Aggregiert werden Peer-Messwerte und Peer-Klassifikationswahrscheinlichkeiten;
          die erzeugten Kontextmerkmale umfassen außerdem Gewichtssumme und
          Anzahl der beitragenden Peers.
  \end{itemize}
  \item \textbf{Modellvarianten}
  \begin{itemize}
    \item Model~1: eigene vier Messwerte, State und aggregierter Kontext
          (12 Eingabemerkmale).
    \item Model~2: eigene Messwerte, State, eigene Klassenwahrscheinlichkeiten
          und aggregierter Kontext (15 Eingabemerkmale).
    \item Vergleich mit Baseline ohne Kontext sowie mit einer zusätzlichen
          transformierten Baseline, bei der eigene und aggregierte Messwerte
          kombiniert werden.
  \end{itemize}
  \item \textbf{Hypothese}
  \begin{itemize}
    \item Kontextinformationen können bei verteilter Aktivität die
          Erkennung verbessern; ob dies Fehlalarme reduziert, ist separat
          anhand einer Falsch-Positiv-Rate zu prüfen.
  \end{itemize}
\end{itemize}

\subsection{Noch zu erklären}
\begin{itemize}
  \item Notebook~05 verwendet alle sieben anderen Knoten; eine räumliche oder
        logische Nachbarschaftstopologie mit weniger Peers ist nicht ersichtlich.
        Klären, ob das Experiment ein vollständiges Mesh oder ein abstrahiertes
        P2P-Modell darstellen soll.
  \item Begründen, warum die Gewichte von der Angriffswahrscheinlichkeit
        abhängen, und die Sigmoid-Gleichung sowie deren Parameter herleiten.
  \item Beschreiben, ob die aggregierten Merkmale in der Praxis über ein
        Netzwerk übertragen werden oder nur synchron aus vorbereiteten
        Datensätzen berechnet werden. Latenz, Paketformat, Ausfälle und
        Vertrauensannahmen sind in den Notebooks nicht experimentell belegt.
  \item Die Begriffe ``Konsens'' und ``lokale Erkennung'' präzisieren:
        Die Baseline wird laut Notebook~04 auf gepoolten Daten trainiert;
        sie ist daher nicht ohne Weiteres ein separat trainiertes Modell
        pro EVSE.
  \item Leakage-Schutz, Vorverarbeitung, Normalisierung und verwendete
        Modellparameter transparent dokumentieren.
\end{itemize}

\subsection{Grafikbedarf}
\begin{itemize}
  \item \textbf{Datenflussdiagramm:} lokale Merkmale
        $\rightarrow$ Random-Forest-Wahrscheinlichkeiten
        $\rightarrow$ gewichtete Peer-Aggregation
        $\rightarrow$ Kontextmodell und Klassenvorhersage.
  \item \textbf{Netzwerkdiagramm:} Node A bis H und die tatsächlich
        verwendeten Peer-Verbindungen; nicht vor einer Klärung der Topologie
        als Ring oder anderes Nachbarschaftsmodell darstellen.
\end{itemize}
